% !TEX program = pdflatex
% Core proof-mechanics document for j15marti/pnp.
% P vs NP remains open.

\documentclass[11pt]{article}
\usepackage[T1]{fontenc}
\usepackage[utf8]{inputenc}
\usepackage[margin=0.82in]{geometry}
\usepackage{amsmath,amssymb,mathtools,amsthm}
\usepackage{booktabs,longtable,tabularx,array}
\usepackage[table]{xcolor}
\usepackage{enumitem}
\usepackage[most]{tcolorbox}
\usepackage{hyperref}
\usepackage{microtype}
\usepackage{fancyhdr}
\usepackage{lastpage}

\definecolor{proved}{HTML}{1B7F3A}
\definecolor{openstatus}{HTML}{B15C00}
\definecolor{retired}{HTML}{A12A2A}
\definecolor{candidate}{HTML}{4B5FA5}
\definecolor{softblue}{HTML}{EEF4FA}
\definecolor{softgreen}{HTML}{EDF8F0}
\definecolor{softorange}{HTML}{FFF4E8}
\definecolor{softred}{HTML}{FCEEEE}

\hypersetup{
 colorlinks=true,
 linkcolor=candidate,
 urlcolor=candidate,
 citecolor=candidate,
 pdftitle={P vs NP Research Program --- Core Proof Mechanics},
 pdfauthor={P vs NP Research Project}
}

\pagestyle{fancy}
\fancyhf{}
\lhead{\small P vs NP Research Program}
\rhead{\small Core Proof Mechanics}
\cfoot{\small \thepage\ /\ \pageref{LastPage}}
\setlength{\headheight}{14pt}

\newcommand{\statusproved}{\textcolor{proved}{\textbf{PROVED/EXACT}}}
\newcommand{\statusopen}{\textcolor{openstatus}{\textbf{OPEN}}}
\newcommand{\statusretired}{\textcolor{retired}{\textbf{FALSIFIED/RETIRED}}}
\newcommand{\statuscandidate}{\textcolor{candidate}{\textbf{CANDIDATE}}}
\newcommand{\same}{\textsc{same\_realization}}
\newcommand{\poly}{\operatorname{poly}}

\newtcolorbox{statusbox}{
 colback=softorange,colframe=openstatus,boxrule=0.8pt,arc=2pt,
 left=7pt,right=7pt,top=6pt,bottom=6pt
}
\newtcolorbox{provedbox}{
 colback=softgreen,colframe=proved,boxrule=0.7pt,arc=2pt,
 left=7pt,right=7pt,top=6pt,bottom=6pt
}
\newtcolorbox{warningbox}{
 colback=softred,colframe=retired,boxrule=0.7pt,arc=2pt,
 left=7pt,right=7pt,top=6pt,bottom=6pt
}
\newtcolorbox{candidatebox}{
 colback=softblue,colframe=candidate,boxrule=0.7pt,arc=2pt,
 left=7pt,right=7pt,top=6pt,bottom=6pt
}

\theoremstyle{plain}
\newtheorem{theorem}{Theorem}[section]
\newtheorem{lemma}[theorem]{Lemma}
\newtheorem{corollary}[theorem]{Corollary}
\theoremstyle{definition}
\newtheorem{definition}[theorem]{Definition}
\theoremstyle{remark}
\newtheorem{remark}[theorem]{Remark}

\title{\bfseries P vs NP Research Program\\[0.2em]
\Large Core Proof Mechanics}
\author{LSPV / SRAS / Exact-3SAT Research Synthesis}
\date{2026-10-05}

\begin{document}
\maketitle

\begin{statusbox}
\textbf{Global status.}
\[
\boxed{\text{$P$ versus $NP$ remains OPEN.}}
\]
This file is a theorem map and proof-mechanics record. It does not claim an
unconditional separation. The central unresolved bridge is
\[
\boxed{
\text{arb. poly.-time exact SAT decider}
\Longrightarrow
\text{forced, source-matched, fully charged reduction structure.}
}
\]
\end{statusbox}

\tableofcontents
\bigskip

\section*{How to read this document}

This file is organized around the logical dependencies of the research
program rather than the chronology of individual experiments. The early
sections isolate the abstract proof mechanics: what counts as a legitimate
reduction, how future-relevant distinctions are tracked, how source and
representation are separated, and where a hypothetical polynomial SAT
algorithm would have to enter the argument. The middle sections record closure
theorems and no-go results that prevent us from charging work that an ordinary
polynomial computation can in fact reuse, reread, recompute, or compress. The
exact-3SAT sections then serve as a concrete laboratory: they identify
structural reductions that really are cheap and therefore must be credited by
any general lower-bound accounting system.

The main distinction throughout is between three kinds of statements:
\emph{accounting identities}, which are exact once the model is fixed;
\emph{structural theorems}, which prove that some representation or relation
must exist; and \emph{cost theorems}, which would show that acquiring or using
that structure is expensive. Most of the remaining P-vs-NP difficulty lies in
the third category, together with the arbitrary-algorithm extraction theorem
needed to connect an unknown SAT decider to the charged structure.

\section{Proof architecture}

The project has two interacting tracks. They should be viewed as complementary
rather than competing approaches. The first track asks what would be required
for an unrestricted lower-bound proof. The second track stress-tests those
requirements on explicit 3-SAT structure, so that the abstract accounting
never charges a resource that a concrete polynomial or exact-exponential
procedure can already eliminate cheaply.

\paragraph{Separation mechanics.}
The Source-Reduction Accounting System (SRAS) and Contractive Source Extraction
(CSE) program ask whether every hypothetical polynomial-time exact SAT decider
must expose enough source-grounded reduction structure to contradict a
representation-independent acquisition barrier.

\paragraph{Exact 3-SAT mechanics.}
The Liu/type-1-chain program classifies concrete cheap structure: local
pruning, direct pair/triple correlations, outside mediators, mediator
elimination, constant-width closure, exact rank-$2/3/4$ resolved relations,
and hypermatching. These results determine what must be credited as cheap
before lower-bound accounting begins.

\begin{candidatebox}
\textbf{Core separation spine.}
\[
\boxed{
\begin{array}{c}
\text{hypothetical polynomial SAT decider }A\\
\Downarrow
\quad\\
\text{source-to-reduction extraction (OPEN)}\\
n-O(\log n)\text{ valid continuation reduction}\\
\Downarrow\quad\\
\text{SRAS accounting}\\
\mathrm{CONTRACT}\ \vee\ \mathrm{UNIVERSAL\ EVALUATOR}\ \vee\ \mathrm{EXTERNAL}\\
\Downarrow\\
\text{close every ordinary branch at superpolynomial cost}\\
\Downarrow\\
P\ne NP.
\end{array}}
\]
The final implication is proved conditionally. The extraction and barrier
hypotheses remain open.
\end{candidatebox}

\section{Quantifier and accounting discipline}

\begin{itemize}[leftmargin=1.4em]
\item Semantic multiplicity is not computational state count.
\item Computational state count is not runtime.
\item Semantic existence is not acquisition.
\item Small representation is not small source.
\item Verification is not construction.
\item DAG reuse is charged once; recomputation is charged again.
\item Static reread is a legal recovery authority unless explicitly forbidden.
\item Never replace $\forall A\,\exists F_A$ by $\exists F\,\forall A$ without proof.
\item Restricted-model lower bounds do not become unrestricted runtime lower
bounds without a model-matched extraction theorem.
\end{itemize}

\begin{candidatebox}
A polynomial SAT algorithm need not look like assignment search, resolution,
a branching tree, an OBDD, or an explicit semantic quotient. Any proof that
assumes one of these representations without extraction is model-restricted.
\end{candidatebox}

\section{Boolean reduction threshold}

This section fixes the amount of reduction that would be needed if a
continuation problem really began with Boolean mass on the order of $2^n$.
It is intentionally only an accounting calibration. Its purpose is to define
the target scale precisely: a polynomial-size residual family requires the
residual exponent to fall from $n$ to $O(\log n)$. The difficult theorem is
not this calculation, but showing that an arbitrary SAT computation is forced
to realize an independently defined reduction of that magnitude.

\begin{theorem}[BRTHRESH1]
If an independently defined reduction of factor $2^x$ leaves residual scale
$n^d$ from an initial scale $2^n$, then
\[
\boxed{x=n-d\log_2 n.}
\]
Equivalently, with $D(n)=n-x(n)$, polynomial residual scale is exactly
\[
\boxed{D(n)=O(\log n).}
\]
\end{theorem}

\begin{proof}
From $2^{n-x}=n^d$, take $\log_2$.
\end{proof}

\begin{corollary}
If $D(n)=\omega(\log n)$, then $2^{D(n)}=n^{\omega(1)}$ is superpolynomial.
\end{corollary}

\begin{remark}
BRTHRESH1 is algebraic accounting. It does not prove that an arbitrary fast
SAT algorithm literally acquires $n-O(\log n)$ assignment-elimination bits.
\end{remark}

\section{Continuation equivalence and safe compression}

A central lesson of the project is that counting raw assignments, histories,
or provenance labels is too coarse. What matters is whether two states can be
merged without changing any future behavior that the declared consumer may
request. Continuation equivalence formalizes exactly that criterion. It gives
a representation lower bound for any summary that must answer all consumers
in $U$, while deliberately stopping short of a runtime claim.

\begin{definition}
Let $\mathcal H$ be histories or residual situations and $U$ a declared set of
legal continuations. Define
\[
h\equiv_U h'
\iff
\forall u\in U,\quad \Phi(h,u)=\Phi(h',u),
\]
and let $\kappa_U(\mathcal H)=|\mathcal H/\!\equiv_U|$.
\end{definition}

\begin{lemma}[CONTSEP1]
If $\sigma:\mathcal H\to Z$ has a deterministic decoder satisfying
$D(\sigma(h),u)=\Phi(h,u)$ for all $h,u$, then
\[
h\not\equiv_U h'\Longrightarrow \sigma(h)\ne\sigma(h').
\]
Hence
\[
\boxed{|\operatorname{range}\sigma|\ge\kappa_U(\mathcal H).}
\]
\end{lemma}

\begin{proof}
Choose a continuation separating $h$ and $h'$. Equal summaries would force
the decoder to return the same answer in both cases.
\end{proof}

\begin{corollary}[NO2K1]
If $|\mathcal H|=2^n$ and a continuation-complete summary has
$|\operatorname{range}\sigma|\le2^{n-k}$, then
\[
\boxed{k\le n-\log_2\kappa_U(\mathcal H).}
\]
\end{corollary}

\begin{warningbox}
This is a representation theorem, not a time lower bound.
\end{warningbox}

\section{Source lineage, provenance, and SRAS}

SRAS is the bookkeeping layer that sits between semantic distinguishability
and machine cost. Its purpose is to prevent information from disappearing
from the proof merely because the computation changes representation. A
future-relevant distinction can be stored explicitly, encoded into another
object, localized to a smaller source, safely quotiented away, reconstructed
later, or supplied externally. The accounting system records which of these
events actually occurred and what ordinary computation was available to
support it.

For a derived object $q$:
\[
\boxed{
J_{\rm live}^{U}(q)
\subseteq J_{\rm sem}(q)
\subseteq J_{\rm syn}(q)
\subseteq J_{\rm origin}(q).
}
\]

Prime labels, hashes, authenticated replay, and \same can establish source
identity, but
\[
\boxed{\text{unique provenance}\not\Rightarrow\text{expensive acquisition}.}
\]

\begin{definition}[SRAS state]
Use
\[
\mathcal A_t=
(X,K,Y,Z,H,J,P,B,U,\Phi_U,V,\Psi_V,R,E,C,Q,G,L,W,\mathsf A,
\Omega_D,\Omega_G,\mathrm{Fate}).
\]
\end{definition}

Useful fate labels are
\[
\mathrm{LIVE},\ \mathrm{ENCODED},\ \mathrm{LOCALIZED},\
\mathrm{QUOTIENTED},\ \mathrm{COMBINED},\
\mathrm{IRRELEVANT},\ \mathrm{TERMINAL},\ \mathrm{EXTERNAL}.
\]

Define valid continuation reduction by
\[
\boxed{
K_{s\to t}
=
\log_2\frac{\kappa_U(S_s)}{\kappa_U(S_t)}.
}
\]

Charge a computation DAG by
\[
\boxed{
C(\mathcal D)=\sum_{v\in V(\mathcal D)}c(v).
}
\]

\begin{candidatebox}
\textbf{SRAS-NFR.}
Every future-relevant unit of valid reduction must remain accounted for as
LOCALIZED, ENCODED/COMBINED, SAFELY QUOTIENTED, REACQUIRABLE, or EXTERNAL.
The cost of these routes remains the load-bearing issue.
\end{candidatebox}

\begin{theorem}[SRAS-ROUTECLOSE1: deterministic causal route completeness]
Fix a deterministic computation interface at time $t$. Suppose all future
behavior can depend only on:
\[
R_t=\text{retained machine state},\qquad
A_t=\text{declared accessible/re-readable source},
\]
and future external inputs $E_{>t}$.
Let $h_0,h_1$ be two legal histories. If
\[
R_t(h_0)=R_t(h_1),\qquad
A_t(h_0)=A_t(h_1),
\]
and the two runs receive the same future external inputs, then every future
deterministic machine state and output is identical.

Therefore, if some legal continuation requires incompatible future responses
for $h_0$ and $h_1$, at least one distinguishing contribution must be:
\[
\boxed{
\text{retained/encoded}
\ \vee\
\text{recoverable from declared accessible source}
\ \vee\
\text{supplied by future external information}.
}
\]
A distinction may instead be QUOTIENTED/IRRELEVANT only after continuation
equivalence proves that no legal future consumer distinguishes it.
\end{theorem}

\begin{proof}
Induct on future deterministic transition steps. Equal current state, equal
accessible source answers, and equal future external inputs give equal next
states; iteration gives equal future outputs. The contrapositive gives the
route-completeness statement.
\end{proof}

\begin{remark}
SRAS-ROUTECLOSE1 closes route completeness only in the declared deterministic
interface model. It proves no branch expensive and does not prove that static
SAT correctness forces $n-O(\log n)$ valid reduction. LOCALIZED, COMBINED, and
other SRAS labels refine the retained/recoverable cases.
\end{remark}

Define
\[
\boxed{
\mathcal C_{F,U,M}(k)
=
\min_{\mathcal D}
\{C(\mathcal D):K_U(\mathcal D)\ge k\}.
}
\]
The desired target is
\[
\boxed{
\mathcal C_{F_n,U_n,M}(n-O(\log n))=n^{\omega(1)}.
}
\]

\section{Extraction and recursive contraction}

This is the current load-bearing section of the separation program. The
execution trace of a polynomial algorithm is already polynomial, so merely
extracting a transcript cannot yield a contradiction. What is needed is a
\emph{structural extraction}: from the behavior of an arbitrary SAT decider,
one must obtain either genuinely smaller recursively consumable source objects
or a globally supported evaluator whose ordinary acquisition/use cost can be
proved superpolynomial. The recurrence theorems below explain why fractional
contraction would be powerful if such an extraction theorem were available.

A polynomial-time algorithm has a polynomial-size ordinary execution trace.
Extracting that trace is not the hard part.

\begin{theorem}[TRACECLOSE1]
If deterministic $A$ runs in polynomial time and $E$ is a polynomial-time
transducer, then
\[
E(x,\operatorname{Tr}_A(x))
\]
is polynomial-time computable.
\end{theorem}

Thus the generic claim that a polynomial trace plus polynomial audit yields a
superpolynomial resource is RETIRED.

\begin{theorem}[EXTCONSUMER1: fixed extension consumer]
Define the solver-independent extension consumer
\[
W_{\rm EXT}(F,\rho)
=
\operatorname{SAT}(F\upharpoonright\rho),
\]
where $\rho$ is any legal partial assignment/restriction of $F$.
For every exact deterministic SAT decider $A$,
\[
E_A(F,\rho)
=
A(F\upharpoonright\rho)
\]
is a correct universal evaluator for $W_{\rm EXT}$, with only the charged cost
of constructing/accessing the restricted instance plus the cost of $A$ on that
instance.
\end{theorem}

\begin{remark}
EXTCONSUMER1 removes the earlier source-definition ambiguity for an
all-restrictions consumer: the continuation semantics are fixed independently
of the solver. It does not give CSE, because ordinary single-variable
restrictions have child scale $N-O(1)$ rather than $N^b$ with $b<1$.
The theorem also relativizes.
\end{remark}

\begin{candidatebox}
\textbf{CSE1: Contractive Source Extraction.}
For every deterministic exact SAT decider $A$ and sufficiently large hard
instance $F$ of scale $N$, extract an effectively acquired source-bound object
with recursively consumable children $Q_1,\dots,Q_m$ satisfying
\[
m\le N^{a+o(1)},\qquad
|J_{\rm live}^{U}(Q_i)|\le N^{b+o(1)},\quad b<1,
\]
and
\[
C_{\rm acquire}+C_{\rm combine}\le N^{c+o(1)}.
\]
The construction must be exact, recursively closed, wrapper-covariant,
source-matched, fully charged, and contain no hidden current-scale SAT call.
\end{candidatebox}

CSE1 remains OPEN.

\begin{theorem}[RECURCLOSE1]
Suppose an $A$-derived SAT procedure performs $n^{r+o(1)}$ local work, makes at
most $n^{\eta+o(1)}$ recursive SAT calls, and every call has size at most
$n^{\theta+o(1)}$ with $0\le\theta<1$. Under recursive closure and no hidden
current-scale call,
\[
\boxed{
T(n)\le n^{C+o(1)},
\qquad
C=\max\left\{r,\frac{\eta}{1-\theta}\right\}.
}
\]
\end{theorem}

\begin{theorem}[MULTICHANNEL1]
For finitely many channels $j$,
\[
\boxed{
C=
\max\left\{
r,\max_j\frac{\eta_j}{1-\theta_j}
\right\}.
}
\]
\end{theorem}

\begin{theorem}[SELECTCHARGE1]
If $n^{u+o(1)}$ selector operations each call a SAT decider of exponent $K$
on residuals of size $n^{\delta+o(1)}$, then the implementation contains cost
\[
\boxed{n^{u+\delta K+o(1)}.}
\]
Current-scale calls $\delta=1$ reintroduce the unknown exponent.
\end{theorem}

\section{Randomness, replay, and universal evaluators}

Earlier versions of the program explored random challenge strings,
SAME\_REALIZATION constraints, and replay resistance as possible ways to make
source dependence computationally binding. The surviving conclusion is more
limited but useful: freshness can prevent advance knowledge, and provenance
can prevent substitution by an unrelated witness, but neither prevents a
compact evaluator from answering correctly after the challenge is revealed.
This section records that closure explicitly so that freshness is not
mistaken for a time lower bound.

\begin{lemma}[Fresh Snapshot]
For IID fair bits $R_1,R_2,\dots$, if $\tau$ is a stopping time for the natural
filtration, then
\[
(R_{\tau+1},\dots,R_{\tau+b})\mid\mathcal F_\tau
\sim U(\{0,1\}^b).
\]
\end{lemma}

\begin{theorem}[FRESH-EVAL1]
If a polynomial-size committed source $R$, a polynomial-size maintainable
history quotient $h$, and a polynomial-time evaluator
$E(R,h,u)=\Phi_R(h,u)$ answer every legal challenge, then polynomially many
fresh/adaptive rounds remain polynomial-time.
\end{theorem}

Hence
\[
\boxed{\text{freshness}+\same\not\Rightarrow\text{expensive response}.}
\]

For static SAT, reread, replay, memoization, preprocessing, and compact
retained summaries all belong to the UNIVERSAL-EVALUATOR branch.

\section{Bounded-width closure no-go}

The clean-core experiments repeatedly produced useful relations that were
syntactically hidden behind mediators or several rounds of local inference.
WIDTHCLOSE1 supplies a guard against overinterpreting such hiding. If all
needed consequences live at constant clause width, then a brute-force
saturation algorithm can enumerate the entire inference space in polynomial
time. Therefore a genuine acquisition barrier must survive constant-width
closure and other compact global polynomial evaluators.

\begin{theorem}[WIDTHCLOSE1]
The number of non-tautological clauses of width at most $w$ is
\[
\boxed{
N_w(n)=\sum_{i=0}^{w}2^i\binom ni.
}
\]
Exhaustive width-$w$ resolution saturation is polynomial in $N_w(n)$. For
constant $w$,
\[
\boxed{N_w(n)=n^{O(w)}=\poly(n).}
\]
\end{theorem}

\begin{corollary}
A useful relation hidden only behind constant-width resolution cannot by
itself supply a superpolynomial SRAS acquisition barrier.
\end{corollary}

\begin{lemma}[GUARDED-UNSAT1]
If $H(W)=\bigwedge_i C_i(W)$ is unsatisfiable and $Q(X)$ is disjoint, then
\[
\boxed{
\bigwedge_i(Q\vee C_i)\equiv Q.
}
\]
\end{lemma}

If $H$ is 2-SAT, affine, bounded-treewidth, or constant-width-resolution easy,
the hidden relation is cheaply recoverable.

\section{Relativization and conditional separation}

The accounting framework by itself can be phrased in a way that relativizes,
so a successful proof needs a later step that uses something specific about
ordinary computation. The present candidate is source transparency: an
ordinary primitive has an explicit implementation that may be expanded and
charged, whereas an arbitrary oracle answer is an opaque external operation.
This observation explains where a nonrelativizing theorem could live, but it
does not yet supply the necessary lower bound. Algebrization remains a
separate barrier.

Baker--Gill--Solovay imply that a fully relativizing proof cannot establish
ordinary $P\ne NP$.

SRAS marks oracle-supplied power by
\[
\Omega_G=\mathrm{EXTERNAL}.
\]

A possible nonrelativizing theorem may inspect ordinary implementation/source
structure and explicitly stop at an opaque oracle node $y=O(x)$. It is invalid
to declare oracle queries expensive by fiat.

The Google Drive source audit recovers a more concrete candidate:
the \emph{Absolute Source-Complexity Separator} (ASCS).  Fix an ordinary
universal machine $U$ and define
\[
\mathrm{Low}_U(c;m,B,\tau)
=
\{y\in\{0,1\}^m:
\exists q,\ |q|\le B,\ U(q,c)\text{ outputs }y\text{ within }\tau\}.
\]
Then
\[
D_{\rm ASCS}(y)=1
\iff
y\notin\mathrm{Low}_U(c;m,B,\tau).
\]
The Drive report proves the exact counting bound
\[
|\mathrm{Low}_U|<2^{B+1},
\]
the corresponding density statement, and a genuine oracle-replacement test:
an oracle can inject strings that are fast for $U^O$ but not in the absolute
ordinary-source set $\mathrm{Low}_U$.  Thus this particular source predicate
really can fail under oracle substitution.

What remains open is the \emph{Fast Source Certificate} (FSC): the ordinary
source-image membership/production relation must admit a bounded
co-nondeterministic verification circuit of the size required by the RMV
reconstruction. Direct replay generally charges the full source production
time.  Therefore ASCS supplies a concrete location for nonrelativization, but
not yet the needed lower-bound or reconstruction theorem.

Algebrization remains a separate open barrier.

\begin{theorem}[SRAS-SEP1: conditional]
Assume an infinite 3-SAT family $F_n$ such that every polynomial exact SAT
decider must account for $n-O(\log n)$ valid continuation reduction. Assume an
exact exhaustive SRAS theorem forces that reduction into
\[
\mathrm{CONTRACT}\ \vee\
\mathrm{UNIVERSAL\ EVALUATOR}\ \vee\
\mathrm{EXTERNAL}.
\]
Suppose CONTRACT yields a model-matched contradiction, every ordinary
UNIVERSAL-EVALUATOR realization carrying the required reduction costs
$n^{\omega(1)}$, and EXTERNAL is unavailable to ordinary computation. Then
\[
\boxed{P\ne NP.}
\]
\end{theorem}

\begin{warningbox}
SRAS-SEP1 is only a conditional implication. Its extraction and barrier
hypotheses are the unsolved content.
\end{warningbox}

\section{Exact 3-SAT structural mechanics}

The exact-3SAT track is used as a hostile control for the abstract proof
mechanics. Whenever a relation, mediator pattern, quotient, or continuation
map can be exposed and exploited efficiently in this concrete setting, SRAS
must credit that operation rather than count it as unexplained hardness. The
results below therefore serve two purposes: they improve the structural
analysis of the Liu/type-1-chain regime, and they progressively narrow the
kind of resource that could plausibly support an unrestricted lower bound.

A positive type-1 chain has state space
\[
D=\{0,1\}^3\setminus\{000\},
\qquad |D|=7,
\]
and Liu characteristic
\[
\boxed{\lambda_1=\frac37.}
\]

The current internal type-1 balance baseline is
\[
\boxed{c_0\approx1.3279235672.}
\]

\subsection{Low-degree mediator elimination}

Let
\[
G_y=
\bigwedge_{i\in P}(A_i\vee y)
\wedge
\bigwedge_{j\in N}(B_j\vee\neg y).
\]

\begin{theorem}[MED-ELIM1]
\[
\boxed{
\exists y\,G_y
\iff
\left(\bigwedge_{i\in P}A_i\right)
\vee
\left(\bigwedge_{j\in N}B_j\right)
}
\]
and equivalently
\[
\boxed{
\exists y\,G_y
\iff
\bigwedge_{i\in P}\bigwedge_{j\in N}(A_i\vee B_j).
}
\]
\end{theorem}

Constant mediator degree therefore admits polynomial local elimination.

\subsection{NP-hardness inside the clean-core schema}

The importance of the next theorem is diagnostic. Low outside-variable
degree and the absence of direct chain-only correlations do not by themselves
make the residual formula easy: the current clean-core syntax can still encode
arbitrary 3-SAT. At the same time, the reduction exposes highly structured
equality/disequality constraints after mediator elimination, showing that
NP-hard expressiveness and exact-algorithm exploitability can coexist.

\begin{theorem}[RESIDUAL-HARD1]
SAT restricted to the project class EM2-CLEAN is NP-complete, where base
clauses are disjoint positive type-1 chains, every outside mediator has degree
2 and both polarities, every mediator hides one binary consistency clause
\[
(Q\vee y)\wedge(Q\vee\neg y),
\]
and there are no non-base chain-only residual clauses. The reduction from
arbitrary 3-SAT is linear.
\end{theorem}

This is a structural NP-completeness result, not an exponential lower bound.

\subsection{Exact consistency characteristics}

Two independent type-1 chains have $\lambda_{\rm ind,2}=9/49$.

\begin{theorem}[CONSISTENCY-CHAR1]
For coordinate equality and disequality,
\[
\boxed{\lambda_{=} = \frac{81}{391}},
\qquad
\boxed{\lambda_{\ne} = \frac{27}{130}}.
\]
\end{theorem}

The relative DLS factors are
\[
\boxed{\frac{391}{441}},
\qquad
\boxed{\frac{130}{147}}.
\]

\subsection{Closed-form rank-2/3/4 characteristic}

Mediator elimination can expose relations involving two, three, or four
type-1 chains. Rather than analyze these cases separately, CLAUSECHAR1 gives
one closed form for the Liu characteristic of every single non-tautological
resolved clause. This is useful because it converts syntactic low-degree
mediation into a quantitative DLS credit that can be compared uniformly
against independent chains.

Let
\[
K_{u,v}=2^{-d(u,v)},\qquad
N=K^{-1},\qquad
z=N\mathbf1,\qquad
t=\mathbf1^\top z=\frac73.
\]
For a local literal set $L$, let $B(L)\subseteq D$ be the states falsifying
all literals in $L$, and define
\[
a(B)=
z_B^\top
\bigl(N[B,B]\bigr)^{-1}
z_B.
\]

\begin{theorem}[CLAUSECHAR1]
Let $C$ be one non-tautological resolved clause supported on $r$ distinct
type-1 chains, with local falsifying sets $B_i$. Then
\[
\boxed{
\lambda(C)=
\left[
\left(\frac73\right)^r
-
\prod_{i=1}^{r}a(B_i)
\right]^{-1}.
}
\]
\end{theorem}

The proof uses the product Hamming kernel and a Schur complement on the
satisfying principal submatrix.

Exhaustive enumeration gives
\[
324+1512+1296=3132
\]
rank-$2/3/4$ resolvent schemas.

The weakest exact characteristics are
\[
\boxed{\lambda_{\min,2}=\frac{81}{440}},
\qquad
\boxed{\lambda_{\min,3}=\frac{729}{9236}},
\qquad
\boxed{\lambda_{\min,4}=\frac{6561}{193856}}.
\]

Relative to independent chains:
\[
\boxed{
g_2=\frac{440}{441},\quad
g_3=\frac{9236}{9261},\quad
g_4=\frac{193856}{194481}.
}
\]

\begin{provedbox}
Every non-tautological degree-2 mediator resolvent gives strict DLS credit.
\end{provedbox}

\section{RESOLVE-HYPERMATCH and concentration}

The resolved-relation hypergraph separates two qualitatively different
regimes. If useful relations are spread across many chains, a disjoint
hypermatching gives multiplicative DLS credit. If they are concentrated on a
small set of chains, hypermatching alone can fail even though every individual
relation is favorable. The latter case is not treated as evidence of hardness;
instead it identifies a high-influence structural region that should be
attacked by conditioning or branching.

After cheap mediator elimination, form a rank-$2/3/4$ hypergraph on clean
chains. Weight relation $C$ by
\[
w(C)=\log\frac{\lambda(C)}{(3/7)^{r(C)}}.
\]

\begin{theorem}[HYPERMATCH-DENSITY1]
If the resolved-relation hypergraph has $M$ hyperedges of rank at most 4 and
every clean chain lies in at most $\Delta$ hyperedges, then every maximal
hypermatching contains at least
\[
\boxed{\frac{M}{4\Delta}}
\]
hyperedges and covers at least
\[
\boxed{\frac{M}{2\Delta}}
\]
chains.
\end{theorem}

\begin{proof}
One selected edge touches at most four vertices and hence intersects at most
$4\Delta$ hyperedges. A maximal greedy selection must remove all $M$.
\end{proof}

Thus if $M\ge\alpha\nu$ and $\Delta\le D$ for constants $\alpha,D>0$, then
\[
\boxed{
f_{\rm covered}\ge\frac{\alpha}{2D}>0,
}
\]
giving a strict deterministic-balance improvement.

\begin{theorem}[HYPERMATCH-NOFORCE1]
Low outside-mediator degree alone does not imply constant hypermatching
coverage.
\end{theorem}

\begin{proof}
Force every resolved relation to contain one fixed clean chain $T_0$. Then
every hyperedge intersects every other, so every hypermatching has size at
most one and asymptotic coverage tends to zero.
\end{proof}

\begin{theorem}[ACTIVE-DENSITY1]
After separating every clean chain that participates in no non-tautological
resolved relation, let $\nu_{\rm act}$ be the number of remaining active
chains and let $M$ be the number of distinct resolved supports. Since every
support contains at most four chains,
\[
\boxed{
M\ge \frac{\nu_{\rm act}}4.
}
\]
\end{theorem}

\begin{proof}
Every active chain belongs to at least one support, so
$\sum_{e\in E}|e|\ge\nu_{\rm act}$. Since $|e|\le4$,
$4M\ge\nu_{\rm act}$.
\end{proof}

Combining ACTIVE-DENSITY1 with HYPERMATCH-DENSITY1 gives, whenever the active
incidence is bounded by $\Delta$,
\[
\boxed{
f_{\rm covered,act}
\ge
\frac1{8\Delta}.
}
\]
Thus the bounded-incidence side already has a positive constant coverage
fraction; the unresolved case is genuinely high active-chain incidence.

The role of the next theorem is to show that concentration has additional
internal structure. Because every resolved clause has rank at most four,
repeated incidence cannot keep concentrating arbitrarily without either
creating many disjoint outside petals or collapsing onto a still smaller
common center. This bounded-depth phenomenon turns the formerly vague
high-incidence case into a finite local branching problem.

\begin{theorem}[LINK-DICHOTOMY1]
Consider distinct non-tautological single-resolvent clauses of rank at most
four, after duplicate clauses are removed. Fix a chain $T$ of incidence $d$.
The local part of such a clause on $T$ has at most 18 possible non-tautological
patterns: six one-literal patterns and twelve two-literal patterns on distinct
coordinates.

Fix an integer $q\ge2$. For $s=4,3,2,1$, define
\[
B_4(q)=1,
\]
\[
B_s(q)
=
1+18(4-s)(q-1)B_{s+1}(q).
\]
If
\[
\boxed{
d>18B_1(q),
}
\]
then there are $q$ incident clauses with the same local pattern on $T$ whose
supports outside $T$ are pairwise disjoint.
\end{theorem}

\begin{proof}
Pigeonhole the $d$ clauses by the at most 18 local patterns on $T$.
For a fixed center of $s$ chains with fixed local patterns, consider the
outside supports, each of size at most $4-s$. If they have a matching of size
$q$, stop. Otherwise a maximal matching has at most $q-1$ edges; the union of
those edges is a vertex cover of size at most $(4-s)(q-1)$. Apart from at most
one empty outside support, every clause meets this cover. Pigeonhole on a
cover vertex and then on its at most 18 local patterns, reducing to the
$s+1$ case. At $s=4$, fixed support and fixed local patterns determine at most
one distinct clause. The displayed recurrence bounds every no-$q$ family.
\end{proof}

\begin{remark}
LINK-DICHOTOMY1 is the bounded-rank sunflower/link-hypergraph mechanism suggested
by the open-step audit. It converts sufficiently high chain incidence into a
large family of disjoint petals after one common local pattern is fixed.
LOCALPATTERN-BRANCH1 now completes the finite one-/two-literal center-pattern
audit for the declared resolved rank-$\le4$ model. What remains is global
integration through the full preprocessing pipeline, not the local finite
pattern calculation.
\end{remark}

\begin{theorem}[LOCALPATTERN-BRANCH1]
In the resolved rank-$\le4$ type-1-chain interaction model, every nonempty
local center pattern on one chain contains one or two literals.  Up to
coordinate permutation, the 18 possibilities reduce to five sign types.

Let $m$ incident clauses share the same local center pattern $Q(T)$ and have
pairwise chain-disjoint outside petals.  Using the conservative petal factor
\[
g=\frac{440}{441},
\]
the sum of child DLS potentials after branching on $Q$ is bounded relative to
the parent type-1 contribution by
\[
R_Q(m)=a_Q+b_Qg^m,
\]
with the exact table
\[
\begin{array}{c|c|c|c}
Q\text{ type} & a_Q & b_Q &
\min\{m:R_Q(m)<1\}\\
\hline
1+ & 16/21 & 5/7 & 484\\
1- & 5/7 & 16/21 & 433\\
2++ & 20/21 & 3/7 & 968\\
2+- & 19/21 & 4/7 & 790\\
2-- & 33/35 & 4/7 & 1015
\end{array}
\]
Therefore the uniform petal threshold is
\[
\boxed{q_*=1015.}
\]
\end{theorem}

For the two-literal true branches, the exact local characteristics are
\[
\lambda_{2++}=\frac9{20},\qquad
\lambda_{2+-}=\frac9{19},\qquad
\lambda_{2--}=\frac5{11}.
\]
The false branches fix the center literals and expose the disjoint petals.

Combining $q_*=1015$ with LINK-DICHOTOMY1 gives
\[
B_4=1,\qquad
B_3=18253,\qquad
B_2=666307513,\qquad
B_1=36484334181829,
\]
and hence the explicit incidence threshold
\[
\boxed{
D_*=18B_1=656718015272922.
}
\]

\begin{theorem}[RESOLVED-INCIDENCE-DICHOTOMY1]
After removing inactive chains in the declared resolved rank-$\le4$ model:

\begin{enumerate}[leftmargin=1.6em]
\item if the maximum active chain incidence satisfies
\[
\Delta\le D_*,
\]
then ACTIVE-DENSITY1 and HYPERMATCH-DENSITY1 force active-chain coverage
\[
\boxed{
f_{\rm covered,act}
\ge
\frac1{8D_*}
=
\frac1{5253744122183376},
}
\]
with strict resolved-relation characteristic credit;

\item if
\[
\Delta>D_*,
\]
then LINK-DICHOTOMY1 supplies $1015$ pairwise-disjoint petals with one common
local center pattern, and LOCALPATTERN-BRANCH1 supplies a strict local
DLS-potential decrease.
\end{enumerate}
\end{theorem}

\begin{remark}
The constants are enormous and are not intended as practical algorithmic
parameters.  Their significance is structural: the formerly open
high-incidence escape is finite and bounded in this resolved model.  Full
integration into the complete deterministic 3-SAT algorithm, and extension to
every broader residual clause form, remain open.
\end{remark}

\begin{provedbox}
\textbf{Resolved-incidence dichotomy, current scope.}
Within the declared resolved rank-$\le4$ type-1-chain model,
\[
\boxed{
\text{active resolved core}
\to
\begin{cases}
\Delta\le D_* \Rightarrow \text{positive hypermatching coverage},\\
\Delta>D_* \Rightarrow \text{favorable high-incidence branch}.
\end{cases}
}
\]
Thus concentration itself is no longer an open local escape in this model.
The remaining exact-algorithm work is to propagate this dichotomy through the
complete preprocessing/residual schema and compute the resulting global
running-time balance.
\end{provedbox}

\section{How exact-3SAT mechanics feed SRAS}

The point of combining these two tracks is methodological. The exact-3SAT
analysis continuously enlarges the list of transformations that are known to
be cheap. Any later SRAS lower bound must therefore be invariant under those
transformations. In effect, the exact algorithm acts as a red team against the
proof architecture: whenever it finds a compact summary, local elimination,
matching credit, or global polynomial evaluator, that mechanism must be
removed from the candidate hardness resource.

A valid lower-bound barrier must credit:
\begin{itemize}[leftmargin=1.4em]
\item local chain pruning and direct pair/triple correlations;
\item low-degree mediator elimination and constant-width resolution;
\item exact rank-$2/3/4$ relation characteristics and DAG reuse;
\item SCC/2-SAT and Gaussian elimination for affine/XOR structure;
\item bounded-treewidth dynamic programming;
\item compact universal response evaluators.
\end{itemize}

The spread-versus-concentrate dichotomy suggests an SRAS analogue: spread-out
structure yields disjoint reduction blocks, while concentration yields a small
high-influence source region and may support LOCALIZE. The arbitrary-algorithm
theorem is still OPEN.


\section{Drive-recovered structural controls and source-separator clues}

The Google Drive corpus contains two families of results that sharpen the
current proof map but must not be conflated with one another.

\subsection{Restricted component localization: q-Horn/expander line}

Manuscript 21A, independently rechecked in CLAUDE-0314, proves component
localization, majority-component lineage, and linear component-backdoor depth
for its declared proper exact-2P2N expander family.  These are genuine
structural localization theorems inside that formula/restriction model.

They do not establish arbitrary-algorithm CSE.  The same corpus explicitly
retains the affine/Tseitin bypass, representation sensitivity, and the missing
SAME\_REALIZATION acquisition/use theorem for an unrestricted SAT decider.
Accordingly these results are used here as hard-family controls and as examples
of what a successful LOCALIZE theorem can look like, not as an unrestricted
runtime lower bound.

\subsection{RMV source structure: SFSL and orbit-spectral recurrence}

The live LSPV/RMV snapshot and independent CLAUDE-0312 verification establish
the Seed-Fixed Surjectivity Lemma and the Orbit-Spectral Recurrence Lemma.
Their joint lesson is two-sided:

\begin{itemize}[leftmargin=1.4em]
\item generic low-degree extractor geometry is too expressive to force a
nontrivial fixed-seed relation; and
\item a particular accepted proof polynomial may have small orbit-spectral
dimension and therefore a field-valued recurrence.
\end{itemize}

The dominant unresolved transfer is Booleanization.  A recurrence over the
odd-prime field used in the RMV construction does not automatically become a
small relation among the exposed Boolean code-coordinate bits.  Thus the
field-level recurrence remains structural evidence rather than the required
Boolean separator.

\subsection{ASCS/FSC: concrete nonrelativization versus verification cost}

The Absolute Source-Complexity Separator gives a concrete ordinary-source
predicate that can fail under oracle replacement.  This is stronger than the
generic statement that source-sensitive proofs are merely nonblackbox.

The missing Fast Source Certificate asks for an exact, source-bound
membership/production certificate whose deterministic verification circuit is
small enough for the RMV reconstruction interface.  Direct replay can cost the
full source running time, while probabilistic or algebraic verification does
not automatically yield the exact co-nondeterministic circuit needed here.

This isolates the current tension:
\[
\boxed{
\begin{array}{c}
\text{easy extensional verification}\Rightarrow\text{often relativizing},\\
\text{absolute ordinary-source semantics}\Rightarrow\text{can nonrelativize},\\
\text{but exact small verification}\Rightarrow\text{OPEN}.
\end{array}}
\]


\section{Retired generic proof routes}

\begin{longtable}{p{0.42\textwidth} p{0.50\textwidth}}
\toprule
\textbf{Retired claim} & \textbf{Reason}\\
\midrule
\endhead
Large continuation multiplicity implies runtime &
Only representation/state distinguishability follows.\\
Large explicit object implies runtime &
Succinct circuits, DAGs, algebraic bases, and recomputation can compress it.\\
Unique provenance implies expensive reacquisition &
Static reread and deterministic rerun defeat this.\\
Fresh random challenge implies expensive response &
A universal evaluator can answer after revelation.\\
Polynomial trace plus audit implies hard resource &
TRACECLOSE1 keeps the result polynomial.\\
Long constant-width inference implies hard discovery &
WIDTHCLOSE1 saturates fixed width in polynomial time.\\
Large OBDD/MDD/rank/treewidth implies unrestricted runtime &
Representation-specific lower bounds do not transfer without extraction.\\
Low outside-mediator degree implies matching density &
HYPERMATCH-NOFORCE1 gives concentrated counterexamples.\\
\bottomrule
\end{longtable}

\section{Barrier ledger and current dependency}

Red-team every load-bearing lemma against relativization, algebrization,
Natural-Proofs conditions, Horn/dual-Horn, 2-SAT, XOR/Tseitin, bounded
treewidth, independent products, padding, renaming, auxiliary variables,
DEAD/COPY/DELAY/INTERP/RECOMP, and maximal semantic/algebraic oracles.

The surviving dependency is
\[
\boxed{
\begin{array}{c}
A\in P\text{ solves 3-SAT}\\
\Downarrow\quad\text{SAT-specific near-total reduction / CSE (OPEN)}\\
n-O(\log n)\text{ valid reduction}\\
\Downarrow\quad\text{SRAS-ROUTECLOSE1 (PROVED route completeness)}\\
\mathrm{RETAIN/LOCALIZE}\vee\mathrm{REACQUIRE}\vee\mathrm{EXTERNAL}\\
\Downarrow\\
\text{ordinary retained/reacquired branch cost (OPEN)}\\
\Downarrow\\
P\ne NP.
\end{array}}
\]


\section{Open-step audit: clues recovered from earlier work}

This audit was intended to distinguish genuinely missing mathematics from
pieces that had already been proved in earlier branches of the project but had
not yet been integrated into the main dependency graph. It does not close the
overall separation. It does, however, move several tasks from vague open
questions to narrower, testable statements.

The 2026-10-05 audit of prior chats and repository reports changes the status
of several open arrows without closing the overall separation.

\paragraph{SRAS route completeness is no longer the main uncertainty.}
VIEWSEP1/OFFPATH1 supplied the separating-view necessity. Combined with
deterministic causality, SRAS-ROUTECLOSE1 gives an exhaustive retained versus
reacquired versus external location theorem. The unresolved issue is cost, not
the existence of another hidden route.

\paragraph{A fixed arbitrary-algorithm consumer exists.}
SOURCESTOP1 correctly rejected solver-invented counterfactual games. The
restriction consumer $W_{\rm EXT}$ avoids that problem because its semantics
are fixed directly by SAT. The remaining failure is quantitative: ordinary
self-restriction shrinks only by $O(1)$ variables rather than to power scale
$N^b$, $b<1$.

\paragraph{High chain incidence has bounded-rank structure.}
The rank-$\le4$ clean-core relation system admits ACTIVE-DENSITY1 and
LINK-DICHOTOMY1. Thus sufficiently high chain incidence forces many clauses
with a common local center pattern and pairwise disjoint petals. This converts
the previous qualitative concentration obstruction into a finite branching
problem over the 18 possible one-/two-literal local center patterns.

\paragraph{Universal-evaluator hardness remains essentially the hard problem.}
BINARYDOF-UNIFORM correctly separates cheap input-specific shortcuts from one
uniform selector/router. However PHABSORB/STRATFLAT/REACQ/TRACECLOSE controls
show that provenance, explicit strategy size, and fixed-level-PH response
search cannot establish the needed superpolynomial cost under $P=NP$.

\paragraph{The final contradiction is downstream.}
RECURCLOSE1 and the model-matched time-space forbidden-region ledgers already
specify how a true fractional return would be converted into a contradiction.
The current research priority is therefore the extraction/return and
ordinary-source cost theorem, not additional endpoint algebra.


\section{Priority queue}

The open-step audit changes the order of attack. The highest-value tasks are
now those that either close a finite structural case outright or sharpen the
arbitrary-algorithm extraction statement. In contrast, additional
provenance-only constructions, explicit strategy-tree growth, or constant-width
gadgets are lower priority because the existing closure theorems already show
why those mechanisms cannot carry the final lower bound by themselves.

\begin{enumerate}[leftmargin=1.6em]
\item \textbf{Integrate HIGHINC-CLOSE1 globally.}
The five local center-pattern types are now exhausted and the worst threshold is
$q_*=1015$, yielding the explicit structural cutoff $D_*$.  The next task is
to propagate RESOLVED-INCIDENCE-DICHOTOMY1 through the full clean-core
preprocessing pipeline and derive one clean global deterministic running-time
bound, while checking every residual clause form admitted by that pipeline.

\item \textbf{Exploit SRAS-ROUTECLOSE1 rather than reproving route completeness.}
The deterministic location of a future distinction is now exhaustive at the
interface level. The load-bearing question is the cost of retained/encoded or
reacquired ordinary information, plus the separate theorem forcing a
near-total reduction obligation from SAT correctness.

\item \textbf{CSE from the fixed extension consumer.}
Use EXTCONSUMER1 as the canonical solver-independent consumer and seek the
stronger dichotomy
\[
\text{fractional source localization}
\ \vee\
\text{globally supported reusable evaluator}.
\]
The missing quantitative step is to force child/source scale $N^b$, $b<1$, or
an independently expensive global evaluator.

\item \textbf{Universal evaluator / uniform coverage barrier.}
Reuse the BINARYDOF-UNIFORM distinction between cheap per-instance shortcuts
and one answer-independent selector/router. Do not infer hardness from
explicit strategy trees, provenance, or fixed-level-PH response relations.

\item \textbf{Relativization and algebrization.}
The plausible nonrelativizing location remains an implementation-dependent
ordinary-source cost theorem that stops at opaque EXTERNAL nodes. No current
result closes algebrization; every algebraic/rank-based candidate remains a
hostile control rather than a final hinge.
\end{enumerate}

\section{Strict result ledger}

\subsection*{PROVED/EXACT}
\begin{itemize}
\item BRTHRESH1; CONTSEP1; NO2K1 in the declared model.
\item SRAS-ROUTECLOSE1 in the declared deterministic interface model.
\item TRACECLOSE1, EXTCONSUMER1, and deterministic reread/replay no-go results.
\item RECURCLOSE1, MULTICHANNEL1, SELECTCHARGE1 as accounting theorems.
\item Fresh Snapshot, FRESH-EVAL1, WIDTHCLOSE1, GUARDED-UNSAT1.
\item MED-ELIM1 and RESIDUAL-HARD1.
\item CONSISTENCY-CHAR1 and CLAUSECHAR1.
\item exhaustive 3132 rank-$2/3/4$ mediator-resolvent classification.
\item HYPERMATCH-DENSITY1, HYPERMATCH-NOFORCE1, and ACTIVE-DENSITY1.
\item LINK-DICHOTOMY1 as a bounded-rank combinatorial concentration theorem.
\item LOCALPATTERN-BRANCH1 and RESOLVED-INCIDENCE-DICHOTOMY1 in the declared
resolved rank-$\le4$ DLS-potential model.
\item ASCS density/miss/counting and its oracle-break test as an ordinary-source
nonrelativization diagnostic; FSC remains open.
\end{itemize}

\subsection*{PROVED AS CONDITIONAL IMPLICATIONS}
\begin{itemize}
\item SRAS-SEP1.
\item recursive exponent closure once a valid shrinker exists.
\item deterministic balance improvement once a positive hypermatching coverage
fraction is guaranteed.
\end{itemize}

\subsection*{CANDIDATE / FINITE CHECK NEXT}
\begin{itemize}
\item global integration of RESOLVED-INCIDENCE-DICHOTOMY1 into the complete
clean-core deterministic algorithm;
\item extension of the resolved rank-$\le4$ theorem to every residual clause
schema generated by the broader preprocessing pipeline;
\item FSC for the ASCS ordinary-source separator.
\end{itemize}

\subsection*{OPEN LOAD-BEARING LEMMAS}
\begin{itemize}
\item arbitrary-algorithm CSE / fractional source-to-reduction extraction;
\item SAT-specific forcing of near-total valid reduction after route
completeness;
\item representation-independent universal-evaluator/acquisition barrier;
\item global integration of the proved local high-incidence recurrence through the
full exact-3SAT preprocessing/residual pipeline;
\item nonrelativizing and nonalgebrizing hinges;
\item final unconditional $P\ne NP$ contradiction.
\end{itemize}

\section{2026-10-06 synchronized boundary update}

This section synchronizes the core documents with the October 6 LEAD/RED
audits and corresponding prior-discussion decisions. These results sharpen
interfaces and retire shortcuts; they do not close the central
arbitrary-decider or acquisition lower-bound arrows.

\subsection{Exact FSC interface: randomness is not an exact coND certificate}

\begin{theorem}[EXACTNESS-GATE1]
Let \(R(c,y)\) be the ordinary-source image-membership predicate required by
FSC, and let \(V(c,y,\pi;r)\) be a probabilistic verifier with perfect
completeness and soundness strictly below one. Acceptance for one fixed random
seed is not an exact existential certificate for \(R\). Exact acceptance
under the same verifier semantics requires a condition of the form
\[
\exists \pi\ \forall r\; V(c,y,\pi;r)=1
\]
or a separate exact derandomization/soundness theorem.
\end{theorem}

\begin{remark}
Enumerating all \(2^{|r|}\) random strings gives a direct exact repair with a
corresponding \(2^{|r|}\) verification factor. Verifier-efficient PCP/IP/IOP
results therefore do not automatically close FSC in the exact size-bounded
co-nondeterministic interface.
\end{remark}

\begin{theorem}[LOCAL-GLOBAL-GATE1]
A transcript showing that a local RMV evaluator queried source coordinates
\(Q=\{(x_j,a_j)\}\) and produced \(y\) proves only consistency with those
local answers. FSC requires the stronger statement that the answers extend to
one globally legal, perfectly accepted, SAME\_REALIZATION canonical source.
\end{theorem}

Define
\[
\operatorname{EXTEND}_{src}(c,Q)=1
\]
iff there exists one canonical perfectly accepted PCPP/LDE source consistent
with every pair in \(Q\). The corresponding small exact certificate problem
\textbf{FSC-EXTEND1} remains \statusopen. In the current accepted-proof
formulation its natural exact condition has existential proof/source choice
followed by universal perfect-acceptance checking; this is a
\(\Sigma_2^P\)-type quantifier gate unless an additional
construction-specific collapse is proved. Interpolation alone is insufficient
because it can forget accepted-proof membership.

\subsection{Static sufficient state and the answer-vector cap}

\begin{theorem}[STATIC-DS-TYPING1]
For fixed source \(x\) and finite legal query family \(Q_x\), any
deterministic precommitted sufficient representation \(S_x\) with evaluator
\(E\) is exactly a static data structure for
\[
f_x(q)=\operatorname{SAT}(C_x\wedge L_q),
\]
and conversely. This is a typing equivalence, not a lower bound.
\end{theorem}

\begin{theorem}[PREPROCESSING-GAP1]
A query lower bound in a model that leaves preprocessing computation free
cannot imply an acquisition-time lower bound for \(S_x\).
\end{theorem}

\begin{theorem}[ANSWER-VECTOR-CAP1]
For every finite \(Q_x\), the complete answer vector
\[
S_x=(f_x(q))_{q\in Q_x}
\]
is an exact representation using \(|Q_x|\) answer bits and constant
random-access lookup once \(q\) is indexed. Therefore algebraic rank,
product-span rank, or source-coordinate dimension cannot by itself force a
larger semantic sufficient state.
\end{theorem}

Consequently the RMV static-data-structure route has the exact dichotomy:
either the SAT-normalized semantic query family is itself too large for the
endpoint space budget, or construction/acquisition of its answer vector must
be charged and lower-bounded. A pure query-probe lower bound cannot establish
the latter. The direct-bridge interpretation \textbf{RMV-DATASTRUCTURE1} is
therefore \statusretired. The fully normalized diagnostic
\textbf{RMV-DATASTRUCTURE2} remains \statuscandidate.

\subsection{Orbit-query deduplication and normalization}

For the declared SU05/Hadamard realization, ORBITBATCH1 proves that for
\(B=\mathbb F_p^{r/2}\) and invertible orbit map \(A\),
\[
\bigcup_{i\in I}\{A^iv:v\in B\}=B
\]
for every nonempty set of orbit positions \(I\). Seed-by-orbit multiplicity
therefore cannot be charged as distinct field-query addresses.

Equally importantly, \(N=p^r\) is a field-domain cardinality, not the
original SAT input length. Any final forbidden-region comparison must supply
the complete normalization
\[
\boxed{
n_{\rm SAT}\longmapsto |y|
\longmapsto (h,r,\Delta,p)
\longmapsto p^r .
}
\]
Directly interpreting contraction in \(p^r\) as SAT-input contraction is
\statusretired. PAIR2COVER1/FULLPAIR6 remain exact fixed-realization
structural rank statements, but ANSWER-VECTOR-CAP1 prevents promoting that
rank alone to a semantic static-space lower bound.

\subsection{Execution provenance inside SRAS}

Let a deterministic algorithm have realized event history
\[
H_t=(e_1,\ldots,e_t)
\]
and let a polynomial-time observer maintain \(P_t=\phi(H_t)\), including
chained-hash and polynomially maintainable historical-feedback variants.

\begin{theorem}[EXECUTION-HASH-CLOSE1]
If \(A\) runs in polynomial time and the provenance transducer has polynomial
state/update cost, then \(P_A(F)\) is polynomial-time computable from
\((A,F)\) by universal simulation.
\end{theorem}

\begin{theorem}[SEMANTIC-NONDETERMINACY1]
The extensional pair \((F,\operatorname{SAT}(F))\) does not in general
determine an implementation-specific path-sensitive fingerprint \(P_A(F)\).
Different correct implementations can have different realized histories.
\end{theorem}

These statements coexist: the fingerprint is genuinely intensional, but is
not an acquisition lower bound.

\begin{theorem}[SRAS-PROV-ROUTING1]
For fixed deterministic \(A\), typed event boundary \(B\), and deterministic
provenance transducer \(\phi\):
\begin{enumerate}[leftmargin=1.5em]
\item \(P_t=\phi(H_t)\) represents realized history rather than SAT semantics;
\item polynomial-time \(A\) makes \(H_t\) and \(P_t\) polynomially
REACQUIRABLE by replay;
\item a fixed-width fingerprint is a QUOTIENTED representation and cannot be
an exact SAME\_REALIZATION identifier because collisions exist;
\item provenance contributes to valid continuation reduction only when an
independently justified provenance-sensitive consumer \(V,\Psi_V\)
distinguishes the histories.
\end{enumerate}
\end{theorem}

The wrapper observer state is not silently charged as solver-retained state:
one must distinguish \(R_t^A\) from \(R_t^{obs}\).

\paragraph{ORDINARY versus EXTERNAL authority.}
SRAS uses \(\Omega_G=\mathrm{ORDINARY}\) for values derived through declared
ordinary machine operations and \(\Omega_G=\mathrm{EXTERNAL}\) for opaque
oracle/external operations. Equal semantic values may therefore have different
lineage. The provenance layer may propagate this authority difference through
later objects; it does not make EXTERNAL information illegal.

\begin{candidatebox}
\textbf{LINEAGE-BIND1 / REALIZATION-BIND1.}
Attach the realized provenance/authority tag to a downstream object \(q\) and
transport it through declared composition rules \(G\). The candidate use is
to detect substitution or lineage laundering of an EXTERNAL or
different-realization object. This becomes load-bearing only if an
independently defined SAT-derived consumer actually requires the ordinary
lineage. Otherwise it is audit metadata.
\end{candidatebox}

TRACE-INSERTION1 and TRACE-INVARIANCE1 remain exact hostile controls: DEAD,
COPY, DELAY, INTERP, and RECOMP can change literal traces without changing SAT
semantics. Thus raw trace complexity/provenance as a generic SAT-hardness hinge
remains \statusretired. The surviving target is a typed, wrapper-covariant
projection/lineage requirement, not literal trace size.

\subsection{Pair-pruning / Mahaney boundary}

The October 6 pair-pruning audit falsifies the generic rule that equality of
coarse incidence or quadratic-rank signatures permits sound OR-frontier
deletion. Mahaney-style pruning needs a membership-preserving reduction or a
directly sound deletion certificate; structural signature collision alone is
insufficient. This falsification is only for the tested/generic surrogate
route and does not rule out an exact source-transcribed construction.

\textbf{PAIR-REDUCTION-FORCING1} remains \statuscandidate: construct a
source-derived, membership-preserving pair map with polynomial-frontier
pruning, fully charged acquisition, wrapper covariance, and explicit
nonrelativizing/nonalgebrizing audits.

\subsection{Synchronized central dependency graph}

After these updates the load-bearing open chain is still
\[
\boxed{
\begin{array}{c}
\text{arbitrary polynomial exact SAT decider }A\\
\Downarrow\\
\text{SAT-CORRECTNESS-FORCING / arbitrary-decider extraction}\\
\Downarrow\\
\text{CSE / SOURCE-TRANSPORT to a source-matched object}\\
\Downarrow\\
\text{near-total valid reduction or a model-matched reusable evaluator}\\
\Downarrow\\
\text{ordinary-source acquisition/use lower bound}\\
\Downarrow\\
\text{exact endpoint contradiction}\\
\Downarrow\\
\text{nonrelativizing and separately nonalgebrizing discharge}.
\end{array}}
\]

The currently necessary unsolved components are:
\begin{itemize}[leftmargin=1.4em]
\item SAT-CORRECTNESS-FORCING / arbitrary-algorithm CSE;
\item SOURCE-TRANSPORT coverage of arbitrary equivalent implementations;
\item a superpolynomial ordinary-source acquisition/use theorem;
\item closure of the compact universal-evaluator escape without restating SAT hardness;
\item FSC-EXTEND1 at the exact required source/circuit interface;
\item a theorem making ORDINARY rather than EXTERNAL lineage load-bearing for
an independently defined SAT obligation;
\item the full \(n_{\rm SAT}\)-normalized RMV/source length map;
\item exact model matching to any Williams/RMV/circuit endpoint;
\item the full nonrelativizing transfer and a separate nonalgebrization audit;
\item natural-proofs-style audit for any final constructive circuit property;
\item full integration of the exact-3SAT preprocessing/residual pipeline.
\end{itemize}

None of STATIC-DS-TYPING1, ANSWER-VECTOR-CAP1, execution provenance,
PAIR2COVER1/FULLPAIR6, or ORBITBATCH1 supplies one of these missing lower
bounds by itself.

\begin{statusbox}
\textbf{Synchronized status.}
The October 6 results substantially sharpen the boundary of the program but
do not prove the central extraction or acquisition theorem.
\[
\boxed{\text{$P$ versus $NP$ remains OPEN.}}
\]
\end{statusbox}


\section{Canonical repository records}

\begin{itemize}
\item \path{research/explorations/source-reduction/2026-10-05_SRAS_CSE1.md}
\item \path{research/explorations/source-reduction/2026-10-05_SRAS_RANDOM_REPLAY1.md}
\item \path{research/explorations/source-reduction/2026-10-05_RESIDUAL_HARD1.md}
\item \path{research/explorations/source-reduction/2026-10-05_CONSISTENCY_PAIR1.md}
\item \path{research/explorations/source-reduction/2026-10-05_WIDTHCLOSE1.md}
\item \path{research/explorations/source-reduction/2026-10-05_RESOLVE_MATCH1.md}
\item \path{research/explorations/source-reduction/2026-10-05_CLAUSECHAR_HYPERMATCH1.md}
\item \path{research/explorations/source-reduction/2026-10-05_HIGHINC_CLOSE1.md}
\item \path{experiments/source_reduction/high_incidence_pattern_table.py}
\item \path{results/source_reduction/2026-10-05_high_incidence_pattern_table.json}
\item \path{research/explorations/random-challenge/2026-10-02_1345_LSPV5_BINARYDOF_UNIFORM1.tex}
\item \path{research/agents/lspv-3/2026-10-04_0755_LSPV-3_SOURCESTOP1.tex}
\item \path{research/agents/lspv-3/2026-10-04_0832_LSPV-3_VIEWSEP1.tex}
\item \path{research/agents/lspv-4/2026-10-04_0246_LSPV-4_ROBUSTGAP1.tex}
\item \path{research/agents/lspv-4/2026-10-02_0449_LSPV-4_TRACECLOSE1.tex}
\item \path{research/agents/lspv-4/2026-10-01_1546_LSPV-4_REACQ1.tex}
\item \path{research/agents/lspv-2/2026-10-04_1604_LSPV-2_RECURCLOSE1.tex}
\item \path{research/agents/lspv-3/2026-10-04_1635_LSPV-3_MULTICHANNEL1.tex}
\item \path{research/agents/lspv-2/2026-10-04_1623_LSPV-2_SELECTCHARGE1.tex}
\end{itemize}

\section{Final statement}

\begin{statusbox}
The exact-3SAT work identifies and credits increasingly sophisticated cheap
structure. SRAS-ROUTECLOSE1 now makes the deterministic location of a
future-relevant distinction exhaustive at the interface level, while
EXTCONSUMER1 supplies a fixed solver-independent all-restrictions consumer.
ACTIVE-DENSITY1, LINK-DICHOTOMY1, and LOCALPATTERN-BRANCH1 close the
bounded-versus-high-incidence split in the declared resolved rank-$\le4$
type-1-chain model. Global integration into the complete exact-3SAT pipeline
remains open. Random/replay and bounded-width results continue to eliminate
false hardness mechanisms, and recursive-closure
theorems show what a genuine contractive extractor would buy.

But the fundamental quantitative theorem remains missing:
\[
\boxed{
\text{arbitrary polynomial SAT computation}
\Longrightarrow
\text{forced source-matched contraction or superpolynomial acquisition}.
}
\]
Until that arrow and its nonrelativizing/nonalgebrizing barrier are proved,
\[
\boxed{\text{$P$ versus $NP$ remains OPEN.}}
\]
\end{statusbox}


\appendix

\section{SRAS explained: what the accounting system is for}

\subsection{Purpose}

The Source-Reduction Accounting System (SRAS) is a bookkeeping framework for
reasoning about a computation that appears to reduce a large family of
possible continuations to a much smaller future-relevant family.

It is deliberately more conservative than an information-counting argument.
SRAS does not infer time from the number of semantic possibilities. Instead it
asks three separate questions:

\begin{enumerate}[leftmargin=1.6em]
\item \textbf{What future distinctions are genuinely required?}
This is defined relative to a declared continuation or consumer universe $U$.

\item \textbf{Where did the information supporting those distinctions come from?}
This is tracked through source lineage, retained representations, external
inputs, reread authority, and recomputation.

\item \textbf{What did the machine actually pay?}
This is tracked on the computation DAG, crediting reuse and charging
recomputation, construction, access, update, and reconstruction separately.
\end{enumerate}

The point of SRAS is to prevent several invalid identifications:
\[
\text{many semantic cases}
\not\equiv
\text{many machine states}
\not\equiv
\text{large runtime},
\]
and
\[
\text{information exists}
\not\equiv
\text{information has been acquired}.
\]

\subsection{The SRAS state tuple}

The current full state notation is
\[
\mathcal A_t=
(X,K,Y,Z,H,J,P,B,U,\Phi_U,V,\Psi_V,R,E,C,Q,G,L,W,\mathsf A,
\Omega_D,\Omega_G,\mathrm{Fate}).
\]

The coordinates are intended as a typed ledger rather than one monolithic
measure.

\begin{longtable}{p{0.12\textwidth} p{0.77\textwidth}}
\toprule
\textbf{Coordinate} & \textbf{Meaning in the current program}\\
\midrule
\endhead
$X$ & Source identity or source object from which the current object descends.\\
$K$ & Stable clause, residual, or canonical object identity used to distinguish
semantic identity from transient representation.\\
$Y,Z$ & Current semantic values, relations, or state descriptors.\\
$H$ & Derivational or audit history actually realized by the computation.\\
$J$ & Causal/source lineage. The important refinement is
$J_{\rm live}^U\subseteq J_{\rm sem}\subseteq J_{\rm syn}\subseteq J_{\rm origin}$.\\
$P$ & Identity/provenance tags, including prime or Godel-style tags when used.
Magnitude of a tag is not itself a hardness measure.\\
$B$ & Boundary, separator, or interface through which the current object is
consumed by the rest of the computation.\\
$U$ & Declared legal continuation/consumer universe.\\
$\Phi_U$ & Semantic behavior of the current object under all declared
continuations in $U$.\\
$V,\Psi_V$ & Provenance-sensitive observations and their required behavior when
the consumer depends on more than extensional semantics.\\
$R$ & Retained representation currently available to the computation.\\
$E$ & Enabled operations, escapes, alternative algorithms, or legal
transformations still available.\\
$C$ & Certificates already established.\\
$Q$ & Reacquisition/reconstruction mechanism.\\
$G$ & Typed combination, reuse, sharing, or composition rule.\\
$L$ & Loss or quotient event: what information has been discarded or merged.\\
$W$ & Recoverability: what discarded information can be reconstructed later.\\
$\mathsf A$ & Access path: how the machine can read or reach the supporting
information.\\
$\Omega_D$ & Data source authority: input, retained memory, random source,
advice, or another declared data channel.\\
$\Omega_G$ & Operation authority: ordinary implementation, subroutine,
oracle/external primitive, or another execution source.\\
Fate & Final disposition of each future-relevant contribution:
LIVE, ENCODED, LOCALIZED, QUOTIENTED, COMBINED, IRRELEVANT, TERMINAL, or EXTERNAL.\\
\bottomrule
\end{longtable}

\subsection{Source lineage}

SRAS distinguishes four nested source sets:
\[
J_{\rm live}^{U}(q)
\subseteq
J_{\rm sem}(q)
\subseteq
J_{\rm syn}(q)
\subseteq
J_{\rm origin}(q).
\]

\begin{itemize}[leftmargin=1.4em]
\item $J_{\rm origin}$ records everything that causally contributed.
\item $J_{\rm syn}$ records what the current syntax depends upon.
\item $J_{\rm sem}$ removes syntactic dependencies that do not affect meaning.
\item $J_{\rm live}^{U}$ removes semantic dependencies that no declared future
consumer can distinguish.
\end{itemize}

This hierarchy is important because large historical source support can
collapse to a small future-live representation without contradiction.

\subsection{Fate labels}

The Fate coordinate answers: what happened to each contribution that was once
future-relevant?

\begin{description}[leftmargin=1.7em,style=nextline]
\item[LIVE] The distinction is still required by some legal future continuation.
\item[ENCODED] The distinction is stored in another representation.
\item[LOCALIZED] The distinction has been isolated to a smaller source region.
\item[QUOTIENTED] Several histories or source states were safely merged because
they are continuation-equivalent.
\item[COMBINED] The information is represented jointly with other information
through a reusable composed object.
\item[IRRELEVANT] It has been proved irrelevant to the declared consumers.
\item[TERMINAL] No future consumer remains.
\item[EXTERNAL] The distinction is supplied by an external authority such as an
oracle or genuinely new information source.
\end{description}

\subsection{Valid reduction}

For a state $S$, let $\kappa_U(S)$ denote the number of continuation-equivalence
classes still distinguished by the declared consumer universe $U$.

The valid reduction from stage $s$ to stage $t$ is
\[
\boxed{
K_{s\to t}
=
\log_2
\frac{\kappa_U(S_s)}{\kappa_U(S_t)}.
}
\]

This is deliberately semantic/consumer-relative. It is not defined from
runtime.

For example, a transformation that reduces the explicit number of assignments
but leaves every continuation class distinguishable may have large syntactic
shrinkage but
\[
K_{s\to t}=0.
\]

Conversely, a safe quotient can have large $K_{s\to t}$ while costing little
if a compact canonical representation computes the quotient cheaply.

\subsection{The No-Free-Reduction invariant}

The project uses the following accounting rule.

\begin{candidatebox}
Every future-relevant reduction contribution must be accounted for as
LOCALIZED, ENCODED/COMBINED, safely QUOTIENTED/IRRELEVANT, REACQUIRABLE through
a charged recovery path, or EXTERNAL.
\end{candidatebox}

SRAS-ROUTECLOSE1 now proves the corresponding route-completeness statement in
the declared deterministic interface model: an incompatible future response
must be distinguished by retained state, declared accessible/re-readable
source, or future external information. The unresolved lower-bound problem is
the \emph{cost} of the ordinary retained/reacquired routes and the separate SAT
theorem forcing a near-total reduction obligation.

\subsection{Acquisition curve}

The acquisition curve turns that question into a quantitative target:
\[
\boxed{
\mathcal C_{F,U,M}(k)
=
\min_{\mathcal D}
\left\{
C(\mathcal D):
K_U(\mathcal D)\ge k
\right\}.
}
\]

Here $M$ is the declared computation model, $\mathcal D$ is a legal
source-transparent computation DAG, and
\[
C(\mathcal D)=\sum_{v\in V(\mathcal D)}c(v)
\]
charges nodes actually computed.

The P-vs-NP-strength target is
\[
\boxed{
\mathcal C_{F_n,U_n,M}
\bigl(n-O(\log n)\bigr)
=
n^{\omega(1)}.
}
\]

A polynomial SAT algorithm would have to contradict this only after a separate
extraction theorem proves that it really must instantiate the corresponding
reduction object.

\subsection{Worked SRAS example: mediator elimination}

Consider a low-degree mediator
\[
(A\vee y)\wedge(B\vee\neg y).
\]

Eliminating $y$ gives the exact resolvent
\[
A\vee B.
\]

In SRAS terms:

\begin{itemize}[leftmargin=1.4em]
\item The original mediator clauses contribute to $J_{\rm origin}$.
\item The mediator variable itself can become IRRELEVANT after exact elimination.
\item The relation $A\vee B$ is ENCODED as a reusable residual constraint.
\item No reacquisition charge is justified merely because the original
mediator variable disappeared.
\item If $A\vee B$ is a constant-width consequence, WIDTHCLOSE1 says its
discovery belongs to a polynomially saturable inference class.
\end{itemize}

This example illustrates the intended discipline:
semantic elimination is credited immediately when an ordinary polynomial
implementation can expose and reuse it.

\section{CSE explained: the missing extraction theorem}

\subsection{Why extraction is the hard arrow}

A polynomial-time algorithm already has a polynomial execution trace. The
difficult task is not to copy that trace. The difficult task is to extract an
object with the exact structural property needed by the contradiction.

The desired arrow is
\[
\boxed{
A
\Longrightarrow
\text{source-matched contractive or expensive future-relevant object}.
}
\]

An extraction theorem must therefore be stronger than trace recovery,
provenance recovery, or state counting.

EXTCONSUMER1 supplies one canonical solver-independent consumer:
\[
W_{\rm EXT}(F,\rho)=\operatorname{SAT}(F\upharpoonright\rho).
\]
Every exact SAT decider automatically evaluates this consumer. Thus the
remaining CSE problem is not defining counterfactual semantics; it is obtaining
a genuine fractional source contraction or proving that the resulting global
extension evaluator has independently expensive ordinary-source acquisition.

\subsection{Typed CSE contract}

A usable Contractive Source Extraction theorem should specify:

\begin{enumerate}[leftmargin=1.6em]
\item the input computation model for $A$;
\item the source object being extracted from the computation;
\item the exact consumer universe $U$;
\item how the extracted object is constructed;
\item its live source footprint;
\item the number and size of recursive children;
\item how child answers combine;
\item construction, access, update, and reconstruction charges;
\item the recursion rule at smaller scales;
\item wrapper covariance under padding, delay, copying, interpolation, and
recomputation;
\item why no hidden current-scale call to $A$ remains.
\end{enumerate}

The current candidate form is
\[
m\le N^{a+o(1)},\qquad
|J_{\rm live}^U(Q_i)|\le N^{b+o(1)},\quad b<1,
\]
with
\[
C_{\rm acquire}+C_{\rm combine}\le N^{c+o(1)}.
\]

If such a source-matched shrinker exists, RECURCLOSE1 converts the recursive
geometry into an explicit exponent.

\subsection{Extraction failure modes}

A proposed extraction lemma fails if any of the following occur:

\begin{itemize}[leftmargin=1.4em]
\item it defines reduction from the observed runtime of $A$;
\item it assumes the algorithm is a search tree or resolution proof;
\item it extracts only an explicit trace whose polynomial size is already known;
\item construction of the extracted object calls SAT again at the current scale;
\item a wrapper changes the extracted measure without changing computational
difficulty;
\item a compact representation or algebraic evaluator services all consumers
without materializing the extracted object;
\item the theorem works equally well with an opaque oracle and therefore cannot
supply the intended nonrelativizing hinge.
\end{itemize}

\section{Continuation equivalence and consumer design}

\subsection{Why the consumer universe matters}

A quotient is safe only relative to a declared future use.

Two histories can be equivalent for the final SAT bit while inequivalent for a
stronger future task such as:

\begin{itemize}[leftmargin=1.4em]
\item returning a satisfying assignment;
\item answering arbitrary residual-extension queries;
\item producing a canonical witness;
\item answering source-bound replay challenges;
\item supporting a later restriction or separator computation.
\end{itemize}

Therefore every continuation argument must state $U$ before counting classes.

\subsection{Safe merge rule}

If $h\not\equiv_Uh'$, then a deterministic correct computation cannot merge
them into one indistinguishable state unless it does at least one of:

\begin{itemize}[leftmargin=1.4em]
\item retains side information;
\item keeps separate states;
\item rereads/reacquires distinguishing source information later;
\item recomputes a sufficient distinguishing object;
\item receives new external information.
\end{itemize}

This is a correctness constraint. It becomes a lower bound only when all legal
ways of satisfying the constraint have independently proved cost.

\section{Provenance, random challenges, and replay}

\subsection{What provenance establishes}

Provenance can establish that two objects came from different source
realizations even when their current semantic outputs coincide.

This is useful for defining \same constraints and for preventing an argument
from silently swapping witnesses or histories.

It does not imply that the provenance-sensitive object is expensive.

\subsection{Replay attack taxonomy}

The current project uses four standard replay attacks.

\begin{description}[leftmargin=1.7em,style=nextline]
\item[Deterministic rerun]
If the needed object is a deterministic polynomial-time function of the static
input, discard-and-rerun is polynomial.

\item[Transcript regeneration]
A polynomial-time functional of a polynomial execution transcript remains
polynomial by TRACECLOSE1.

\item[Collapse/selection replacement]
Under the counterfactual $P=NP$, polynomially balanced fixed-PH witness or
provenance relations can often be selected in polynomial time.

\item[Compositional quotient]
If a long challenge/restriction history composes into a polynomial-size
canonical quotient, retain the quotient rather than the explicit history.
\end{description}

\subsection{Freshness}

A post-stopping-time random challenge is genuinely unpredictable from the
prior causal state. This blocks challenge-specific precomputation.

It does not block a reusable evaluator
\[
E(R,h,u)
\]
that computes the correct answer after $u$ is revealed.

Hence the current random/replay target is a universal-response-evaluator
barrier, not merely a provenance or freshness theorem.

\section{Clean-core and type-1-chain glossary}

\begin{description}[leftmargin=1.9em,style=nextline]
\item[Type-1 chain]
A single positive 3-clause with local satisfying domain
$D=\{0,1\}^3\setminus\{000\}$ and characteristic $3/7$.

\item[Bare chain]
A type-1 chain before additional residual interactions are incorporated.

\item[Direct correlation]
A residual clause supported entirely on variables belonging to two or more
clean chains.

\item[Outside mediator]
A variable outside the chain triples that appears in clauses coupling chain
literals.

\item[Externally mediated clean core]
A set of bare chains with no direct chain-only rank-2/3 residual coupling,
while remaining interactions pass through outside variables.

\item[Degree-2 mediator]
An outside variable with one positive and one negative occurrence. Exact
elimination yields a single resolvent.

\item[Resolved relation]
The finite-domain relation on the participating chain states after mediator
elimination.

\item[Rank]
The number of distinct chain variables participating in the resolved relation.
A degree-2 mediator produces rank at most 4.

\item[RESOLVE-MATCH]
Collect exposed rank-2 relations, combine constraints on the same chain pair,
compute their exact characteristics, and take a weighted matching.

\item[RESOLVE-HYPERMATCH]
Generalize to rank-$2/3/4$ relations and greedily select disjoint favorable
blocks.

\item[Concentration obstruction]
Many favorable relations can all share the same chain, destroying large
hypermatching coverage even though each relation is individually useful.

\item[High-chain incidence]
The exact-algorithm regime in which one chain participates in many resolved
relations.

\item[Link hypergraph]
For a fixed high-incidence chain $T$, remove $T$ from the supports of its
incident resolved clauses. LINK-DICHOTOMY1 shows that bounded rank forces
either many pairwise disjoint petals or repeated concentration onto a larger
fixed center.

\item[Petal]
The portion of a resolved clause support outside a fixed common center.
Disjoint petals are useful because their DLS gains can potentially multiply
without chain overlap.

\item[Finite center-pattern audit]
LINK-DICHOTOMY1 reduces the high-incidence problem to the 18 possible
one-/two-literal local center patterns, which collapse by symmetry to five sign
types. LOCALPATTERN-BRANCH1 verifies all five. The worst conservative petal
threshold is $q_*=1015$, attained by the two-negative-literal center type.
\end{description}

\section{Characteristic and DLS formulas}

\subsection{Local kernel}

For the seven-state domain $D$, define
\[
K_{u,v}=2^{-d(u,v)},\qquad N=K^{-1},
\]
\[
z=N\mathbf1,\qquad
t=\mathbf1^\top z=\frac73.
\]

The type-1 characteristic is
\[
\lambda_1=t^{-1}=\frac37.
\]

For a local falsifying set $B\subseteq D$, define
\[
a(B)=
z_B^\top
\bigl(N[B,B]\bigr)^{-1}
z_B.
\]

\subsection{Single resolved clause}

For a resolved clause $C$ touching $r$ chains with local falsifying sets
$B_1,\dots,B_r$:
\[
\boxed{
\lambda(C)=
\left[
\left(\frac73\right)^r
-
\prod_{i=1}^{r}a(B_i)
\right]^{-1}.
}
\]

The current exhaustive worst cases are:

\begin{center}
\begin{tabular}{cccc}
\toprule
Rank & Independent $\lambda$ & Worst resolved $\lambda$ & Relative DLS factor\\
\midrule
2 & $9/49$ & $81/440$ & $440/441$\\
3 & $27/343$ & $729/9236$ & $9236/9261$\\
4 & $81/2401$ & $6561/193856$ & $193856/194481$\\
\bottomrule
\end{tabular}
\end{center}

For equality/disequality-correlated chain pairs:
\[
\lambda_{=} = \frac{81}{391},
\qquad
\lambda_{\ne}=\frac{27}{130}.
\]

\subsection{How the characteristic enters DLS}

If a disjoint block has characteristic $\lambda_R$, then the local DLS term
contains $\lambda_R^{-1}$ instead of multiplying the independent single-chain
terms separately.

A favorable relation has
\[
\lambda_R >
\left(\frac37\right)^r.
\]

Its relative DLS multiplier is
\[
g(R)=
\frac{(3/7)^r}{\lambda_R}<1.
\]

This explains why hypermatching is useful: disjoint blocks can be credited
without double-counting chain variables.

\section{The spread-versus-concentration dichotomy}

After all cheap degree-2 mediators are eliminated, let $H_{\rm res}$ be the
rank-$\le4$ hypergraph of non-tautological resolved relations.

There are two broad regimes.

\subsection{Spread-out regime}

If there are linearly many relations and bounded clean-chain incidence, then
HYPERMATCH-DENSITY1 guarantees a constant covered fraction:
\[
M\ge\alpha\nu,\qquad
\Delta\le D
\quad\Longrightarrow\quad
f_{\rm covered}\ge\frac{\alpha}{2D}.
\]

Because every selected resolved relation has positive DLS credit, this regime
improves the internal deterministic type-1 balance.

\subsection{Concentrated regime}

If incidence is unbounded, a small set of chains may intersect most exposed
relations. Hypermatching can then be small.

This is not evidence of hardness. It is a new structural signal:
a small source region has high influence over many local constraints.

The next exact-algorithm target is to prove that conditioning or branching on
such a high-incidence chain creates a favorable recurrence.

The corresponding SRAS intuition is that concentration may support LOCALIZE,
but no arbitrary-algorithm localization theorem has been proved.


\section{New-agent intensional and separator boundary results}

\subsection{Restricted extensional access can be eliminated}

\textbf{RESTRICTED-ACCESS-ELIMINATION1 (PROVED/EXACT in the stated deterministic compiler model).}
Let \(H\) be a SAT family and let a deterministic polynomial-time compiler
\(C\) see a candidate exact SAT procedure only through returned decision bits.
If every adaptive query issued by \(C\) lies in \(H\), and a fixed ordinary
polynomial-time algorithm \(B\) correctly decides every instance in \(H\),
then \(B\) can replace the exact SAT procedure and reproduce the complete
adaptive transcript and output of \(C\).  The proof is induction on the
adaptive query number.

This is a necessity theorem, not arbitrary-algorithm extraction.  It
relativizes and continues to hold when both simulations receive the same
algebraic-oracle resources.  It therefore supplies neither the
nonrelativizing nor the nonalgebrizing hinge.

\subsection{Global intensional gate and source transport}

\textbf{GLOBAL-INTENSIONAL-GATE1 (CANDIDATE).}
The current architecture asks a successful SAT-native compiler to combine:
global use of SAT correctness beyond an independently easy restricted family;
intensional dependence on ordinary source/implementation information; explicit
charging of source acquisition, inspection, transformation, verification, and
reuse; and a landing object that gives genuine contraction or belongs to a
representation class with an unconditional superpolynomial lower bound.
No theorem currently forces all four properties for an arbitrary polynomial
SAT decider.

\textbf{SOURCE-TRANSPORT1 (CANDIDATE).}
Free minimization to a canonical semantically equivalent source is not an
acceptable repair for wrapper dependence.  Efficient syntactic
canonicalization need not remove arbitrary semantics-preserving wrappers,
whereas fully extensional canonicalization becomes a program/circuit
minimization or equivalence problem.  The proposed replacement fixes a typed
source language and allowed wrapper monoid and requires explicit
polynomial-time certificate transport with bounded description and
verification overhead.  Coverage of arbitrary relevant implementations
remains open.

\subsection{Causal random provenance}

An external IID random bitstream can force causal ordering and prevent
challenge-specific precomputation, but the existing reacquisition, randomized
closure, transcript closure, strategy-flattening, and PH-provenance controls
do not imply superpolynomial post-revelation work.

Accordingly, causal random provenance remains
\textbf{CANDIDATE/DIAGNOSTIC ONLY}.  The strengthened open target is an
unavoidable post-revelation-work theorem that survives polynomial evaluators
and selectors, memoization, replay, ordinary rerun, and fixed-level-PH
response relations.

\subsection{Exact affine-separator boundary}

\textbf{LINEAR-SEPARATOR-EQUIV1 (PROVED/EXACT in the stated finite Boolean
range model).}
For nonempty \(R\subseteq\mathbb F_2^m\), a nonconstant affine Boolean
function \(D(x)=a\cdot x+b\), \(a\ne0\), satisfying \(D(r)=0\) for every
\(r\in R\) exists if and only if
\[
\operatorname{aff}(R)\ne\mathbb F_2^m.
\]
Such a separator is exactly half-dense.  If \(0\in R\), then \(b=0\), and the
criterion is equivalently
\[
\operatorname{span}(R)\ne\mathbb F_2^m.
\]

Thus succinct constraint detection for an upstream linear code does not by
itself establish a common parity separator for the exact final Boolean RMV
range.  After nonlinear/adaptive extraction and Boolean postprocessing, one
must separately prove final-range affine-rank deficiency.  That obligation
remains \textbf{OPEN}.  LINEAR-SEPARATOR-EQUIV1 is finite linear algebra,
relativizes and algebrizes, and is not a runtime lower bound.

\subsection{Current strict ledger}

The new results do not alter the global status.  \(P\) versus \(NP\) remains
open.  Arbitrary-algorithm SAT-native extraction/CSE, FSC, an
ordinary-source acquisition lower bound, exact final RMV affine-rank
deficiency, unavoidable post-revelation work, and the required
nonrelativizing/nonalgebrizing transfer remain open.

Canonical reports:
\begin{itemize}[leftmargin=1.4em]
\item \path{research/2026-10-05_1818_LEAD-0001_global_intensional_gate1.md}
\item \path{research/2026-10-05_1900_PROOF-0001_restricted_access_elimination.md}
\item \path{research/agents/RED/2026-10-05_RED-0003_canonical_source_minimization_attack.md}
\item \path{research/2026-10-05_2201_PROOF-0004_linear_separator_boundary.md}
\end{itemize}

\section{Barrier and hostile-control checklist}

Before promoting a candidate lemma to the proof spine, apply the following
audit.

\subsection{Representation controls}

\begin{itemize}[leftmargin=1.4em]
\item Can a DAG share the purported repeated work?
\item Can memoization or preprocessing eliminate reacquisition?
\item Can a compact circuit, algebraic basis, or decision procedure represent
the entire consumer map?
\item Does explicit tree/MDD/OBDD growth disappear in another representation?
\end{itemize}

\subsection{Easy-family controls}

Test the claimed resource on:

\begin{itemize}[leftmargin=1.4em]
\item Horn and dual-Horn SAT;
\item 2-SAT/SCC;
\item affine/XOR/Tseitin with Gaussian elimination;
\item bounded-treewidth formulas;
\item independent products.
\end{itemize}

A resource that is large on one of these families but whose decision problem
is polynomial is not, by itself, an unrestricted runtime measure.

\subsection{Wrapper controls}

Test DEAD, COPY, DELAY, INTERP, RECOMP, padding, variable renaming, and
auxiliary-variable encodings.

The measured quantity should not inflate merely because the same computation
is represented with more dead syntax, copied source support, delay steps, or
alternative encodings.

\subsection{Oracle controls}

Replace semantic subcomputations by a maximal oracle.

Ask exactly which step fails.

If the proof continues unchanged, it is likely still relativizing.

If it fails only because the oracle is declared expensive, the argument is not
valid.

A genuine nonrelativizing step must use ordinary implementation/source
structure that is absent from the opaque oracle interface.

\subsection{Algebrization control}

Even a nonrelativizing source argument may continue to work when the machine is
given low-degree algebraic extensions of its oracle. This requires a separate
audit.

\section{Status-label conventions}

This document uses the following meanings.

\begin{description}[leftmargin=1.8em,style=nextline]
\item[PROVED/EXACT]
A mathematical statement proved for the explicitly declared model or a finite
computation exhaustively verified in the declared finite domain.

\item[COMPUTATIONAL]
An experimental finite result, fit, search, or measurement. It is not an
asymptotic theorem.

\item[CONDITIONAL]
The implication is proved if explicitly stated hypotheses hold, but at least
one hypothesis is not established.

\item[CANDIDATE]
A proposed lemma, invariant, or architecture whose load-bearing content remains
unproved.

\item[FALSIFIED/RETIRED]
A proposed generic hinge has a counterexample, closure theorem, representation
escape, or other reason it cannot support the intended unrestricted
conclusion.
\end{description}

\section{Compact glossary of recurring project terms}

\begin{longtable}{p{0.23\textwidth} p{0.67\textwidth}}
\toprule
\textbf{Term} & \textbf{Meaning}\\
\midrule
\endhead
SRAS & Source-Reduction Accounting System: the typed ledger for semantic
reduction, source lineage, representation, reuse, loss, recovery, access, and
external authority.\\
CSE & Contractive Source Extraction: the desired theorem extracting a
recursively smaller, fully charged, source-matched object from an arbitrary SAT
decider.\\
\same & Requirement that a response or witness belong to the same committed
source realization rather than an arbitrary semantically valid replacement.\\
Continuation equivalence & Equality of all required future behavior under the
declared continuation universe $U$.\\
Universal evaluator & Compact reusable algorithm/representation that can answer
all legal future challenges after they are revealed.\\
Reacquisition & Reconstruction of discarded information by reread, rerun,
replay, recomputation, or another charged recovery path.\\
Clean core & Residual family of disjoint type-1 chains whose remaining
cross-chain interactions are externally mediated rather than direct chain-only
clauses.\\
Resolved relation & Finite-domain constraint obtained after eliminating a
mediator.\\
Hypermatching & Collection of pairwise chain-disjoint resolved relations so
their DLS credit can be used without overlap.\\
Concentration & Regime where many exposed relations share a small number of
chains, defeating large hypermatching coverage.\\
\bottomrule
\end{longtable}


\begin{thebibliography}{9}
\bibitem{Liu2018}
S. Cliff Liu,
\emph{Chain, Generalization of Covering Code, and Deterministic Algorithm for $k$-SAT},
ICALP 2018 / arXiv:1804.07901.

\bibitem{BGS1975}
T. Baker, J. Gill, R. Solovay,
\emph{Relativizations of the $P=?NP$ Question},
SIAM Journal on Computing 4(4), 1975.

\bibitem{AW2009}
S. Aaronson, A. Wigderson,
\emph{Algebrization: A New Barrier in Complexity Theory},
ACM Transactions on Computation Theory, 2009.

\bibitem{RR1997}
A. Razborov, S. Rudich,
\emph{Natural Proofs},
Journal of Computer and System Sciences 55(1), 1997.

\bibitem{BW2001}
E. Ben-Sasson, A. Wigderson,
\emph{Short Proofs are Narrow---Resolution Made Simple},
Journal of the ACM 48(2), 2001.

\bibitem{IPZ2001}
R. Impagliazzo, R. Paturi, F. Zane,
\emph{Which Problems Have Strongly Exponential Complexity?},
Journal of Computer and System Sciences 63(4), 2001.

\bibitem{Tovey1984}
C. A. Tovey,
\emph{A Simplified NP-Complete Satisfiability Problem},
Discrete Applied Mathematics 8(1), 1984.
\end{thebibliography}

\end{document}
